% =============================================================
% §2 Background
% =============================================================
\section{Background}
\label{sec:background}

\subsection{LLM Inference: Prefill, Decode, KV Cache}
\label{sec:bg-inference}

\begin{figure*}[t]
  \centering
  \includegraphics[width=0.70\textwidth]{llm_inference_lifecycle.pdf}
  \caption{An LLM inference request's two-phase lifecycle.
    \emph{Prefill} runs once (compute-bound) over the whole
    prompt, emitting the first token and building the KV cache;
    \emph{Decode} then runs $L_o$ sequential, memory-bandwidth-bound
    steps, each reading and extending the KV cache. The KV cache
    is pinned to the GPU that ran Prefill, so dispatching Decode
    elsewhere forces KV migration or a Prefill re-run---the core
    locality tension a cross-server scheduler must navigate.}
  \label{fig:lifecycle}
\end{figure*}

Modern LLMs are decoder-only
Transformers~\cite{vaswani2017attention,brown2020gpt3} whose
inference has two phases with markedly different hardware
characteristics
(Fig.~\ref{fig:lifecycle}). \textbf{Prefill}
runs once per request: it ingests the prompt of $L_p$ tokens
in a single compute-saturating forward pass, producing the
first output token and
a \emph{key-value cache} (KV cache) for every layer.
\textbf{Decode} then runs
iteratively, once per output token ($\sim$10--80~ms each on
A100), reusing and extending the KV cache; the sequence of
$L_o$ steps dominates
end-to-end latency for long generations.

The KV cache is large (0.5--2.5~MB per token per layer;
$\sim$800~MB for a 1000-token LLaMA-13B conversation) and
\textbf{bound to the GPU where Prefill
executed}: serving Decode elsewhere requires either copying
hundreds of MB across the network or re-running Prefill.

Beyond Transformer LLMs, AIGC workloads also include
text-to-image diffusion models whose
inference invokes a shared U-Net at every denoising
step---structurally different but
sharing the model-residency and batching characteristics that
motivate our work.

\subsection{Continuous Batching and Modern Serving Systems}
\label{sec:bg-batching}

\textbf{Continuous (iteration-level) batching}---introduced by
Orca~\cite{yu2022orca}, refined by
vLLM~\cite{kwon2023vllm} (PagedAttention) and
Sarathi-Serve~\cite{agrawal2024sarathi} (chunked
prefill)---recomposes the batch every decoder iteration:
finished requests are evicted, queued ones join mid-flight.
Per-request batching overhead is $<5\%$ in decode
(memory-bandwidth-bound) but $\sim$30\% in prefill
(compute-bound)---the phase-dependent factors our simulator
adopts (\S\ref{sec:simulator}). These engine-level
optimizations operate
\textbf{within one server}; the upstream question of
\emph{which server} should serve each request is the focus of
this paper.

\subsection{Cloud-Edge Collaborative LLM Deployment}
\label{sec:bg-cloudedge}

Production LLM services increasingly span heterogeneous compute
pools: one or more cloud servers with high-end accelerators
(A100/H100, $\ge$80~GB of GPU memory, VRAM) capable of hosting
LLaMA-70B-INT8
($\sim$70~GB) at high throughput; a pool of edge servers
(data-center A100 down through workstation T4 and embedded
Jetson Orin) with 16--64~GB VRAM, 10--50 TFLOPS, and
30--250~W thermal design power (TDP) that fit smaller
LLaMA-7B/13B variants; and a
network fabric with 30--65~ms cloud--edge and 8--23~ms
edge--edge round-trip latency. The paradigm is driven by
\textbf{latency} (edge proximity reduces time-to-first-token,
TTFT),
\textbf{cost} (edge GPUs are 1--10$\times$ cheaper per FLOP and
2--10$\times$ lower power), and \textbf{data sovereignty}.
Phase-disaggregation and migration systems
(\S\ref{sec:related}) redesign the deployment itself; we
address the complementary \emph{initial placement} question
they leave open: given a fixed heterogeneous pool, where
should each request go?

\subsection{Scheduling Approaches}
\label{sec:bg-sched}

Cloud-edge scheduling has historically used heuristic
schedulers (HEFT~\cite{topcuoglu2002heft},
Round-Robin/Least-Loaded/Shortest-Queue), metaheuristics
(GA~\cite{holland1975adaptation},
PSO~\cite{kennedy1995pso}), and learning-based schedulers
(Decima~\cite{mao2019decima},
RLScheduler~\cite{zhang2020rlscheduler}, and---closest in
topology---A3C-R2N2~\cite{tuli2022a3cr2n}, an asynchronous
advantage actor-critic with a gated-recurrent-unit (GRU) +
residual encoder for generic edge-cloud tasks). None natively
encodes AIGC physics: heuristics treat tasks as independent
black boxes, metaheuristic fitness functions require
hand-engineered AIGC terms---precisely what we automate---and
existing learned schedulers predate the LLM era. We survey
this landscape in depth in \S\ref{sec:related} and benchmark
all three families in \S\ref{sec:eval}.
